\chapter{Un nouveau fonds : les manuscrits mathématiques de Leibniz}\label{ch:leibniz}

Ce livre est né d'une demande : reprendre la méthode sur les manuscrits mathématiques de Leibniz, dont une partie est
présentée sur le site du projet PhiLiUMM :
\url{https://eman-archives.org/philiumm/mathematical-manuscripts}. Ce chapitre applique
au fonds la grille du chapitre~\ref{ch:grille} et en tire un plan de travail. Les constats
datent du 3 novembre 2026. Ils ont été faits depuis un terminal, sur quelques pages et
quelques manifestes, pas sur l'ensemble du fonds : on les refait avant de s'appuyer
dessus.

\section{Deux sites, deux rôles}

Il faut distinguer d'abord ce que montre chaque site, parce que le site de PhiLiUMM
n'est pas celui qui détient le fonds.

\begin{description}
  \item[Le projet PhiLiUMM] (\emph{The Philosophy of Leibniz in Light of his Unpublished
        Mathematical Manuscripts}), financé par le Conseil européen de la recherche
        (ERC, ADG~101020985, 2021--2026) et rattaché au laboratoire SPHERE, travaille avec
        le Leibniz-Archiv de Hanovre. Sa page d'équipe nomme en premier David Rabouin
        (CNRS). Le site, sur la plateforme EMAN, présente une sélection de manuscrits en
        neuf rubriques (construction des courbes et des équations, perspective et géométrie
        projective, fondements de la géométrie, \emph{Fundamenta calculi}, justification
        du calcul différentiel, dyadique, angles de contact, \emph{ars combinatoria},
        caténaire). Chaque manuscrit a une présentation et une transcription en HTML, le
        latin de Leibniz avec les formules en notation LaTeX. Les images ne sont pas sur
        ce site : il renvoie à celui de la bibliothèque.
  \item[La Gottfried Wilhelm Leibniz Bibliothek] (GWLB, Hanovre) détient les papiers. Ses
        collections numérisées (\url{https://digitale-sammlungen.gwlb.de}, sur le logiciel
        Kitodo.Presentation) rangent les manuscrits mathématiques sous la cote LH~35
        (\emph{Leibniz-Handschriften zur Mathematik}). Une recherche sur « LH 35 » y renvoie
        311 notices. Chaque unité a une cote fine (par exemple LH~35, 3~A~35) et un
        manifeste IIIF.
\end{description}

\section{La grille, appliquée}

\paragraph{Le fonds est-il servi entier ?} Pour les unités LH~35 qui sont numérisées, oui :
le manifeste d'une unité liste toutes ses vues. Reste à établir quelle part de LH~35 est
numérisée. Le site de PhiLiUMM parle d'environ 17\,000 feuillets mathématiques de Leibniz,
dont moins de la moitié a été transcrite. Le premier travail est donc un inventaire, sur le
modèle de celui de Gallica (chapitre~\ref{ch:gallica}) : une ligne par unité, avec la cote,
le nombre de vues, l'adresse du manifeste, la date du relevé.

\paragraph{Comment les images arrivent-elles ?} Par le protocole IIIF, et c'est la bonne
nouvelle. Les manifestes sont au format IIIF Presentation~2, à une adresse stable de la
forme \fichier{https://digitale-sammlungen.gwlb.de/content/<numéro>/manifest.json}. Chaque
vue a un service IIIF Image (niveau~1) qui sert l'image à la largeur demandée. Les
manifestes et les images répondent avec l'en-tête \texttt{Access-Control-Allow-Origin: *},
ce qui veut dire que n'importe quelle page web peut les lire dans le navigateur. C'est la
situation de Gallica, pas celle de Montpellier : pas besoin de relais, et aucune image n'a
à être stockée. Les vues relevées font autour de 4\,800 pixels de large.

\begin{terminal}
curl -s https://digitale-sammlungen.gwlb.de/content/00067960/manifest.json \
  | python3 -c "import json,sys; m=json.load(sys.stdin); print(m['label'], m['license'])"
\end{terminal}

\paragraph{Que permet la bibliothèque ?} Les deux manifestes relevés, dont un de LH~35,
portent dans leur champ \texttt{license} la \emph{Public Domain Mark}~1.0 de Creative
Commons, avec la mention « No Copyright - Public Domain Marked » et l'attribution à la GWLB.
La bibliothèque publie aussi une \emph{Open Digitisation Policy}, en PDF, qu'on n'a pas
pu lire depuis le terminal (le texte n'en est pas extractible) : il faut la lire, et
vérifier que la marque figure sur chaque manifeste qu'on utilise. Leibniz est mort en
1716, ses manuscrits sont dans le domaine public ; ce sont les images, les notices et les
transcriptions d'autrui qui demandent qu'on regarde les conditions.

\paragraph{Les transcriptions existantes.} C'est la vraie difficulté de ce fonds, et elle
est double.

D'abord, il y a une édition. La série~VII de l'édition de l'Académie (\emph{Sämtliche
Schriften und Briefe}, Reihe~VII, \emph{Mathematische Schriften}) publie les écrits
mathématiques par volumes ; le site de PhiLiUMM la cite (par exemple « A~VII-6 »). Chaque
unité retenue doit donc être cherchée dans l'édition avant d'être transcrite, comme on le
fait pour le fonds Grothendieck avec les éditions de Maltsiniotis. Une unité déjà éditée
n'est pas à refaire ; au mieux, on l'affiche à côté de son fac-similé en renvoyant à
l'édition, si ses conditions le permettent. On n'a pas vérifié lesquels des volumes sont
en ligne ni sous quelle licence.

Ensuite, il y a le travail de PhiLiUMM lui-même. Ses transcriptions sont sur son site, sans
licence qu'on ait trouvée. Il en prépare un encodage en TEI, qu'il présente comme
expérimental : transcriptions d'abord saisies en ODT avec les formules en LaTeX,
converties en XML par l'outil Odette, puis reprises dans d'autres éditeurs. Il a aussi
entraîné des modèles de reconnaissance d'écriture (HTR) sur l'écriture de Leibniz :
détection des zones (RF-DETR), des lignes (Kraken), et un modèle de reconnaissance affiné
à partir de FoNDUE-GD, avec un taux d'erreur par caractère de 8,33\,\% annoncé. Les modèles
et les jeux d'entraînement sont déposés sur Zenodo et Hugging Face. Une équipe a déjà, en
somme, fait le travail le plus dur de la chaîne. On ne le refait pas dans son coin : on lui
écrit.

\paragraph{La main et la langue.} Latin du \textsc{xvii}\ieme{} siècle, abréviations,
notations de Leibniz qui changent d'une année à l'autre (ses signes d'égalité, ses
\emph{omn.}, ses premiers $\int$), ratures et ajouts entre les lignes, figures dans le
texte. Les règles de transcription du fonds Grothendieck s'appliquent (les deux registres,
le doute marqué, la variante gardée), mais la notation d'époque pose une question qu'un
fonds du \textsc{xx}\ieme{} siècle ne pose pas : faut-il transcrire le signe de Leibniz ou
son équivalent moderne ? La réponse de ce livre est de garder le signe dans la
transcription et de ne moderniser que dans la lecture modernisée, qui est un autre
fichier.

\paragraph{Les liens profonds.} Les liens du site de PhiLiUMM vers le visualiseur de la
GWLB portent des paramètres de session (\texttt{cHash}). Une variante de l'un d'eux, sans
ces paramètres, répond par une erreur~404. On ne fonde donc rien sur ces liens : on cite la
cote et l'adresse du manifeste, qui ne changent pas.

\section{Le plan de travail}

\begin{enumerate}
  \item \textbf{Écrire avant de commencer}, à l'équipe de PhiLiUMM et au Leibniz-Archiv,
        pour présenter le projet et demander ce qu'ils accepteraient : partir de leurs
        transcriptions, réutiliser leurs modèles HTR, citer leurs numéros. Une réponse
        négative change le plan ; une absence de réponse ne vaut pas accord.
  \item \textbf{L'inventaire} : les notices LH~35 de la GWLB, une ligne par unité, avec la
        cote, le nombre de vues, le manifeste et sa licence, relevés par un script qui
        attend entre deux requêtes et se nomme dans son \emph{user agent}. On y ajoute pour
        chaque unité sa présence dans la série~VII et sur le site de PhiLiUMM.
  \item \textbf{Le choix d'un premier lot} : quelques unités non éditées, courtes, d'un
        sujet que l'on connaît. Vingt vues, pas plus, pour éprouver la chaîne
        entière : transcription, relecture contre le fac-similé, statut, page de lecture.
  \item \textbf{Le site} sur le modèle de germain-project : le fac-similé lu en direct par
        IIIF à côté de la transcription, aucun fichier image dans le dépôt, un
        \fichier{.gitignore} qui refuse \fichier{*.jpg} et \fichier{*.pdf}.
  \item \textbf{La procédure écrite pour l'assistant} (\fichier{SKILL.md}), adaptée de
        \texttt{transcribe-grothendieck} : latin, abréviations, notation de Leibniz,
        renvoi à l'édition de l'Académie quand l'unité y figure.
  \item \textbf{Lean ensuite, et seulement là où il apporte quelque chose}. Les manuscrits de
        Leibniz contiennent des énoncés arithmétiques et combinatoires (la dyadique, l'art
        combinatoire, des sommes de séries) qu'on peut énoncer et vérifier ; on les énonce
        dans la lecture modernisée, jamais dans la transcription, et l'on cite la
        littérature avant de parler de nouveauté. Le calcul différentiel de Leibniz, lui,
        se lit dans l'histoire des mathématiques et pas dans mathlib.
  \item \textbf{Ne rien publier qui reproduise une transcription d'autrui} sans l'accord de
        son auteur, et garder dans un dépôt privé ce qui attend un accord.
\end{enumerate}

\begin{noteblock}
Ce qui manque à une transcription faite par une machine, c'est la lecture des écritures
anciennes, la description des documents, l'édition critique. Le projet gagne à ce que la
relecture contre le fac-similé soit le travail d'une personne qui sait lire la main, et que
l'assistant fasse le reste.
\end{noteblock}
